% Source: https://artofproblemsolving.com/community/c3h1747731_discarded_questions
% Stable identity: aops-topic-1747731-q04
\textbf{Euler line.} Place the origin at $O$ and let $\mathbf a,\mathbf b,\mathbf c$ be the position vectors of $A,B,C$; then $|\mathbf a|=|\mathbf b|=|\mathbf c|=R$, the circumradius. Let $\mathbf h=\mathbf a+\mathbf b+\mathbf c$. Then
\[
(\mathbf h-\mathbf a)\cdot(\mathbf b-\mathbf c)=(\mathbf b+\mathbf c)\cdot(\mathbf b-\mathbf c)=|\mathbf b|^2-|\mathbf c|^2=0,
\]
so the point $\mathbf h$ lies on the altitude from $A$; in the same way it lies on the altitudes from $B$ and $C$, so it is the orthocentre $H$. The centroid is $G=\dfrac{\mathbf a+\mathbf b+\mathbf c}{3}=\dfrac{\mathbf h}{3}$. Hence $O$, $G$, $H$ are collinear and
\[
OH=3\,OG .
\]

\begin{asy}
usepackage("amsmath");
size(5.5cm);
defaultpen(fontsize(10pt));
pair A=(1,3.6),B=(0,0),C=(6,0);
real d=2*(A.x*(B.y-C.y)+B.x*(C.y-A.y)+C.x*(A.y-B.y));
pair O=((dot(A,A)*(B.y-C.y)+dot(B,B)*(C.y-A.y)+dot(C,C)*(A.y-B.y))/d,(dot(A,A)*(C.x-B.x)+dot(B,B)*(A.x-C.x)+dot(C,C)*(B.x-A.x))/d);
pair G=(A+B+C)/3,H=A+B+C-2*O;
draw(A--B--C--cycle,linewidth(0.9));
draw(circle(O,abs(A-O)),gray);
draw(H+0.35*(H-O)--O+0.7*(O-H),red+dashed);
dot(O);dot(G);dot(H);
label("$A$",A,N);label("$B$",B,SW);label("$C$",C,SE);
label("$O$",O,N);label("$G$",G,S);label("$H$",H,N);
label("\mbox{}",(min(currentpicture,user=true)+max(currentpicture,user=true))/2,black); // keeps the text after the figure black when LaTeX typesets the labels inline
\end{asy}

Therefore
\[
OH=\boxed{\dfrac{3(x^2-16)}{x^2-9}}.
\]
(For this to be a distance, $x$ must satisfy $|x|\ge4$ or $|x|<3$.)
